\documentclass[11pt,a4paper]{article}

\usepackage[utf8]{inputenc}
\usepackage[T1]{fontenc}
\usepackage[brazil]{babel}
\usepackage{geometry}
\geometry{a4paper, margin=1.8cm, top=1.8cm, bottom=1.8cm}
\setlength{\headheight}{15pt}

\usepackage{amsmath,amssymb}
\usepackage{booktabs}
\usepackage{tabularx}
\usepackage{xcolor}
\usepackage{listings}
\usepackage{hyperref}
\usepackage{fancyhdr}
\usepackage{float}
\usepackage{microtype}
\usepackage[skins,breakable]{tcolorbox}

% --- PALETA DE CORES INSTITUCIONAL ---
\definecolor{azulPrimario}{RGB}{18, 52, 102}      % Azul marinho técnico
\definecolor{azulSecundario}{RGB}{35, 85, 145}   % Azul institucional
\definecolor{verdeSucesso}{RGB}{30, 110, 60}     % Verde sóbrio
\definecolor{cinzaFundo}{RGB}{246, 248, 250}     % Cinza muito claro
\definecolor{cinzaBorda}{RGB}{210, 218, 226}     % Cinza sutil
\definecolor{azulCaixa}{RGB}{238, 244, 252}      % Azul pastel suave
\definecolor{laranjaCaixa}{RGB}{253, 246, 237}   % Laranja suave para avisos

\hypersetup{
    colorlinks=true,
    linkcolor=azulPrimario,
    urlcolor=azulPrimario,
    citecolor=azulSecundario
}

% --- ESTILO DE CAIXA DE ANALOGIA DIDÁTICA ---
\newtcolorbox{analogia}[1][]{
    enhanced,
    colback=azulCaixa,
    colframe=azulSecundario,
    coltitle=white,
    fonttitle=\bfseries\small\color{white},
    arc=2.5mm,
    boxrule=0.8pt,
    left=9pt,
    right=9pt,
    top=7pt,
    bottom=7pt,
    title={Analogia Prática: #1},
    breakable
}

\newtcolorbox{conceito}[1][]{
    enhanced,
    colback=cinzaFundo,
    colframe=cinzaBorda,
    coltitle=azulPrimario,
    fonttitle=\bfseries\small\color{azulPrimario},
    arc=2mm,
    boxrule=0.6pt,
    left=8pt,
    right=8pt,
    top=6pt,
    bottom=6pt,
    title={Conceito Técnico: #1},
    breakable
}

\pagestyle{fancy}
\fancyhf{}
\lhead{\footnotesize \textbf{Sistemas Distribuídos} — Evolução Arquitetural do Sistema Hospitalar}
\rhead{\footnotesize UFMA / Engenharia da Computação}
\cfoot{\thepage}
\renewcommand{\headrulewidth}{0.4pt}

\begin{document}

% --- CABEÇALHO FORMAL ---
\begin{center}
    {\large \textsc{Universidade Federal do Maranhão}}\\[1pt]
    {\small Centro de Ciências Exatas e Tecnologia — Departamento de Engenharia da Computação}\\[2pt]
    {\small Disciplina: Sistemas Distribuídos — Docente: Prof. Dr. Luiz Henrique Neves Rodrigues}\\[8pt]
    
    {\Large \textbf{Relatório Técnico de Evolução Arquitetural}}\\[2pt]
    {\large \textit{Modernização com Middleware MOM, Virtualização Isolada em 6 VMs e Análise de Escalabilidade}}\\[5pt]
    
    \textbf{Discentes:} Eduardo dos Santos Oliveira, Fernando da Silva Costa, Herick Vinicius Pinheiro Da Conceição, Samara Santos Viegas e Vyktor Gabryel Costa de Souza\\[6pt]
    \rule{\linewidth}{1pt}
\end{center}

\vspace{-0.3cm}

% =========================================================================
\section{Introdução e Contexto das Alterações}
% =========================================================================

O monitoramento contínuo de sinais vitais em Unidades de Terapia Intensiva (UTI) requer alta disponibilidade, processamento em tempo real e integridade absoluta dos dados médicos. O \textbf{Sistema Distribuído para Auxílio de Dados Hospitalares} foi concebido para coletar, enfileirar, balancear e persistir fluxos massivos de telemetria médica provenientes de sensores clínicos.

Na versão inicial do projeto, o sistema operava com comunicação síncrona direta via protocolo HTTP (Cliente $\to$ Balanceador $\to$ Servidores) e agrupava múltiplos servidores dentro de uma mesma Máquina Virtual compartilhada. Para mitigar riscos inerentes a essa abordagem, as alterações consolidadas concentraram-se em três eixos de evolução:
\begin{enumerate}
    \item \textbf{Adoção de Middleware Orientado a Mensagens (MOM):} Incorporação do RabbitMQ para amortecimento de fluxo, desacoplamento temporal e garantia de entrega incondicional;
    \item \textbf{Segregação Física em 6 Máquinas Virtuais (VMs):} Isolamento de cada componente em nós virtuais autônomos com endereçamento de rede independente, eliminando o compartilhamento de recursos;
    \item \textbf{Atualização da Camada de Observabilidade Gráfica:} Ajuste da topologia visual para mapear em tempo real as 6 máquinas e os novos enlaces de mensageria.
\end{enumerate}

% =========================================================================
\section{Camada de Middleware Orientado a Mensagens (MOM)}
% =========================================================================

\subsection{Conceito de Middleware e o Papel do MOM}

Em computação distribuída, o termo \textbf{Middleware} designa uma camada de software intermediária posicionada entre o sistema operacional e as aplicações finais. Sua função é viabilizar a comunicação e o gerenciamento de dados entre nós heterogêneos de maneira transparente.

Um \textbf{MOM (\textit{Message-Oriented Middleware})} baseia-se na troca assíncrona de dados por meio de mensagens contidas em filas, diferentemente do modelo síncrono convencional (chamada de procedimento ou requisição HTTP bloqueante), no qual o emissor precisa aguardar o receptor concluir o processamento antes de prosseguir.

\begin{analogia}[O Pronto-Socorro e a Sala de Triagem]
Imagine um pronto-socorro onde cada ambulância que chega com um paciente crítico precisasse chamar o médico na porta da sala de cirurgia e aguardar o médico terminar a cirurgia para só então entregar a maca. Se o médico estiver ocupado, a ambulância fica presa no estacionamento, impedida de socorrer outros pacientes na cidade.

O \textbf{MOM} atua exatamente como a \textbf{sala de triagem com protocolo seguro}: a ambulância (Produtor) entrega a ficha clínica na recepção da triagem (RabbitMQ) e é liberada em fração de segundos. A recepção organiza as fichas em ordem de chegada numa fila segura. Assim que um médico (Consumidor) conclui um procedimento, ele retira a próxima ficha disponível e realiza o atendimento. Nenhuma ficha se perde e nenhuma ambulância fica retida na porta do hospital.
\end{analogia}

\subsection{Protocolo AMQP 0-9-1 e Componentes Implementados}

O middleware adota o padrão aberto \textbf{AMQP 0-9-1 (\textit{Advanced Message Queuing Protocol})}, que estabelece as regras formais de serialização, confirmação e roteamento de mensagens. No projeto, a arquitetura foi dividida no módulo \texttt{middleware/}:

\begin{itemize}
    \item \textbf{Produtor AMQP (\texttt{middleware/produtor.py}):} Lê registros contínuos do conjunto de dados clínico (\textit{dataset} CSV) e envia cada leitura como uma mensagem JSON à central de roteamento (\textit{Exchange}) com a chave de roteamento \texttt{dados.paciente}.
    \item \textbf{Broker RabbitMQ (Servidor de Mensageria):} Servidor intermediário encarregado de receber as mensagens, aplicar as regras de roteamento e armazená-las de forma persistente na fila \texttt{fila\_dados\_pacientes}.
    \item \textbf{Consumidor AMQP (\texttt{middleware/consumidor.py}):} Processo de escuta contínua que drena os registros da fila e os encaminha para processamento posterior no balanceador de carga.
\end{itemize}

\begin{analogia}[Exchange, Routing Key e Fila como Central dos Correios]
Para compreender a dinâmica interna do RabbitMQ:
\begin{itemize}
    \item \textbf{Exchange (Central de Distribuição):} É a esteira rolante dos correios que recebe todas as cartas que chegam à agência. Ela não guarda as cartas; apenas olha para a etiqueta e decide para onde despachar.
    \item \textbf{Routing Key (CEP ou Etiqueta de Destino):} O código impresso no envelope (no projeto, \texttt{dados.paciente}) indicando o assunto ou destino da correspondência.
    \item \textbf{Queue / Fila (Caixa Postal ou Escaninho):} A gaveta física onde as cartas ficam guardadas de forma segura e ordenada até que o destinatário venha retirá-las.
\end{itemize}
\end{analogia}

\subsection{Políticas de Confiabilidade e Tolerância a Falhas}

Para garantir que nenhum dado clínico seja extraviado por instabilidade de rede ou pane de energia, foram aplicadas quatro diretrizes estritas:

\begin{enumerate}
    \item \textbf{Persistência de Mensagens (\texttt{delivery\_mode=2}):} Cada mensagem recebida pelo broker é gravada imediatamente em disco rígido (armazenamento não-volátil), garantindo que ela sobreviva mesmo que o RabbitMQ seja reiniciado abruptamente.
    \begin{analogia}[Livro Ata vs. Bilhete Adesivo]
    Em vez de colar um bilhete de papel adesivo na parede que pode cair com uma rajada de vento ou apagar-se no escuro (mensagem apenas em memória RAM), cada prontuário é transcrito com tinta indelével em um livro de registro em cofre de aço (persistência em disco).
    \end{analogia}

    \item \textbf{Confirmação Explícita de Recebimento (\textit{Manual Acknowledgement} — \texttt{basic\_ack}):} O consumidor não descarta a mensagem da fila ao lê-la. O RabbitMQ mantém a mensagem sob custódia e só a remove em definitivo após receber uma confirmação formal (\texttt{basic\_ack}) indicando que o servidor backend processou e salvou os dados com sucesso (código HTTP 200).

    \item \textbf{Rejeição com Reenfileiramento (\textit{Negative Acknowledgement} — \texttt{basic\_nack}):} Caso ocorra timeout, queda de conexão ou indisponibilidade temporária de backend, o consumidor emite a instrução de rejeição com a flag de retenção (\texttt{requeue=True}). O broker reinsere a mensagem no topo da fila para nova tentativa imediata por outro nó disponível.
    \begin{analogia}[Assinatura do Recibo de Entrega]
    É o equivalente a assinar o canhoto do Sedex apenas após abrir a embalagem e conferir que o frasco do medicamento chegou intacto. Se a caixa estiver avariada, o recepcionista recusa a assinatura (\textit{NACK}) e a encomenda retorna para o almoxarifado para reenvio sem prejuízo.
    \end{analogia}

    \item \textbf{Controle de Fluxo e Prevenção de Sobrecarga (\texttt{prefetch\_count=1}):} Instrução que limita o broker a entregar apenas uma mensagem por vez a cada consumidor ativo. Novas mensagens só são despachadas após o processamento da mensagem corrente.
    \begin{analogia}[Uma Pasta Médica por Vez na Mesa]
    Impede que um único profissional retire 50 prontuários de uma vez da triagem e acumule uma pilha sobrecarregada em sua mesa enquanto outros colegas de plantão permanecem ociosos. Cada um atende um paciente por vez, promovendo divisão justa e evitando esgotamento de memória (\textit{backpressure}).
    \end{analogia}
\end{enumerate}

% =========================================================================
\section{Virtualização e Topologia Segregada em 6 Máquinas Virtuais}
% =========================================================================

\subsection{Do Compartilhamento de Recursos ao Isolamento Real de Falhas}

Na primeira versão da infraestrutura, os três servidores de aplicação operavam como processos distintos dentro de uma única Máquina Virtual denominada \texttt{vm-servidores}. Sob os preceitos de Sistemas Distribuídos, esse modelo configurava um falso isolamento: embora fossem portas lógicas distintas, qualquer exaustão de memória RAM, saturação de CPU ou pane no kernel da máquina virtual derrubaria todos os três servidores simultaneamente.

A reformulação implementada no arquivo \texttt{Vagrantfile} segregou o sistema em \textbf{6 Máquinas Virtuais totalmente autônomas}, cada qual executando seu próprio sistema operacional Ubuntu, pilha TCP/IP e alocação dedicada de memória e processamento.

\begin{analogia}[Prédio de Apartamentos vs. Casas Independentes]
Antes, os três servidores eram como três pessoas dividindo o mesmo quarto de uma casa: se o cano daquele cômodo estourasse ou houvesse um curto-circuito na fiação (pane do SO), todos os moradores ficavam desamparados. 

Com as 6 VMs, criamos um condomínio fechado com 6 casas muradas independentes. Se a casa de um servidor sofrer uma pane isolada, as outras 5 casas continuam com energia elétrica, abastecimento e portas abertas para atendimento normal.
\end{analogia}

\subsection{Topologia de Rede Privada e Tabela de Endereçamento}

Todas as máquinas virtuais comunicam-se através de uma rede interna privada (\textit{Host-Only Network} no segmento \texttt{192.168.56.0/24}). Esta configuração garante isolamento contra ruídos externos e simula fielmente a rede local cabeada de um complexo hospitalar.

\begin{table}[H]
\centering
\small
\begin{tabularx}{\textwidth}{l l l X}
\toprule
\textbf{Máquina Virtual} & \textbf{Endereço IP} & \textbf{Porta(s)} & \textbf{Papel Arquitetural e Serviços Alocados} \\
\midrule
\texttt{vm-cliente} & \texttt{192.168.56.10} & — & Produtor de telemetria médica (\texttt{produtor.py}) \\
\texttt{vm-middleware} & \texttt{192.168.56.11} & 5672, 15672 & Broker RabbitMQ (Protocolo AMQP e Painel de Gestão) \\
\texttt{vm-balanceador} & \texttt{192.168.56.12} & 8000 & Consumidor AMQP e Load Balancer (Round Robin) \\
\texttt{vm-servidor1} & \texttt{192.168.56.13} & 8001 & Servidor Backend 01 (FastAPI) e Banco \texttt{server1.db} \\
\texttt{vm-servidor2} & \texttt{192.168.56.14} & 8002 & Servidor Backend 02 (FastAPI) e Banco \texttt{server2.db} \\
\texttt{vm-servidor3} & \texttt{192.168.56.15} & 8003 & Servidor Backend 03 (FastAPI) e Banco \texttt{server3.db} \\
\bottomrule
\end{tabularx}
\caption{Mapeamento topológico das 6 Máquinas Virtuais independentes do projeto.}
\label{tab:vms_topologia}
\end{table}

\subsection{Distinção entre Infraestrutura como Código (IaC) e Execução Local}

Para fins de avaliação do projeto, é fundamental compreender a coexistência harmoniosa de dois ambientes complementares:
\begin{itemize}
    \item \textbf{Ambiente Formal de Infraestrutura como Código (\texttt{Vagrantfile}):} A especificação das 6 máquinas virtuais com provisionamento automatizado é integralmente funcional. Ao invocar o comando \texttt{vagrant up}, o hipervisor VirtualBox inicializa os 6 nós virtuais com adaptadores de rede físicos e serviços independentes.
    \item \textbf{Ambiente de Simulação Ágil (\texttt{executar.sh}):} Para permitir validações rápidas de lógica sem exigir 6 GB de memória RAM dedicada e virtualização aninhada na estação de trabalho do desenvolvedor, o script \texttt{executar.sh} orquestra os mesmos componentes como processos isolados no sistema operacional do host, vinculados às portas de rede locais 8000 a 8003 e comunicando-se com a instância local do RabbitMQ.
\end{itemize}

% =========================================================================
\section{Camada de Apresentação e Visualização em Tempo Real}
% =========================================================================

\subsection{Observabilidade em Sistemas Distribuídos e Protocolo SSE}

Em sistemas distribuídos, a observabilidade consiste na capacidade de inferir o estado interno de um nó a partir das suas saídas externas. Para viabilizar a visualização didática do fluxo de telemetria médica, a interface gráfica (\texttt{visualizacao/index.html}) utiliza \textbf{Server-Sent Events (SSE)}, um protocolo baseado em HTTP que mantém uma conexão unidirecional persistente aberta entre o servidor e o navegador.

\begin{analogia}[O Monitor Cardíaco Contínuo vs. O Questionamento Frequente]
Se o navegador precisasse fazer consultas periódicas de segundo em segundo para saber se há novos dados (\textit{Polling}), seria o mesmo que uma enfermeira entrar a cada dez segundos no quarto do paciente gritando "qual é a pressão arterial agora?". Isso sobrecarrega a rede e desperdiça tempo. 

Com o protocolo \textbf{SSE}, o navegador comporta-se como um \textbf{monitor cardíaco moderno com cabo conectado}: a linha de transmissão permanece aberta e, no milissegundo exato em que um batimento cardíaco ou anomalia clínica é computada, o sinal é enviado instantaneamente para a tela do médico.
\end{analogia}

\subsection{Reformulação da Topologia Gráfica Vetorial (SVG)}

O painel vetorial foi integralmente adaptado para refletir a nova realidade de 6 máquinas:
\begin{itemize}
    \item \textbf{Fronteiras Virtuais Explícitas:} Cada uma das 6 VMs é demarcada por um retângulo tracejado estilizado, contendo seu respectivo nome de host (\texttt{vm-cliente}, \texttt{vm-middleware}, etc.) e endereço IP estático dedicado;
    \item \textbf{Destaque do Middleware:} A exchange, a fila durável e a porta AMQP 5672 ganharam representação visual proeminente, evidenciando o acúmulo temporário e o dreno das mensagens;
    \item \textbf{Animação Física de Pacotes (\textit{Flying Packets}):} Círculos luminosos percorrem o diagrama simulando os pacotes físicos de dados transitando em tempo real: Produtor $\to$ Broker RabbitMQ $\to$ Balanceador $\to$ Servidor Backend selecionado.
\end{itemize}

% =========================================================================
\section{Análise Crítica e Conceitual de Escalabilidade}
% =========================================================================

\subsection{Escalabilidade Horizontal versus Vertical}

A capacidade de um sistema em lidar com o aumento de demanda define sua escalabilidade:
\begin{itemize}
    \item \textbf{Escalabilidade Vertical (\textit{Scale-Up}):} Adicionar mais recursos a uma única máquina (mais processadores, mais memória RAM). Possui barreiras físicas estritas e custo que cresce exponencialmente.
    \item \textbf{Escalabilidade Horizontal (\textit{Scale-Out}):} Adicionar mais máquinas trabalhando em conjunto. É a base dos sistemas distribuídos modernos e a arquitetura adotada em nosso projeto.
\end{itemize}

\begin{analogia}[Trocar o Motor do Veículo vs. Ampliar a Frota]
A escalabilidade vertical é equivalente a trocar o motor de uma ambulância comum pelo motor de um jato supersônico. Além de ter um limite físico intransponível, o custo financeiro é exorbitante e, se aquele veículo único pifar, o hospital inteiro fica sem transporte.

A escalabilidade horizontal equivale a adquirir uma frota com 5 novas ambulâncias de fábrica. A capacidade de atendimento multiplica-se por cinco e, caso uma ambulância precise de manutenção, as outras quatro continuam operando sem qualquer interrupção.
\end{analogia}

\subsection{A Escalabilidade do Projeto é Automática?}

\textbf{Não. A escalabilidade do sistema é horizontal manual (configurável), e não automática (\textit{Auto-Scaling}).}

\begin{analogia}[O Gerente da Agência Bancária vs. A Catraca Inteligente]
Pense em uma agência bancária operando com 3 guichês de atendimento. Se a fila de espera começar a dobrar a esquina, o gerente da agência pode autorizar a abertura de um 4º guichê para diluir o tempo de espera. 

Isso é \textbf{escalabilidade horizontal manual}: o sistema possui a flexibilidade estrutural para comportar novos caixas sem precisar quebrar as regras de atendimento do banco, mas a abertura depende de uma decisão deliberada do gerente. 

A \textbf{escalabilidade automática (\textit{Auto-Scaling})}, por outro lado, seria o equivalente a uma catraca equipada com inteligência computacional que cronometrasse o tamanho da fila e, autonomamente, convocasse novos atendentes e abrisse novos guichês sem qualquer toque humano.
\end{analogia}

\subsection{Capacidades Atuais e Requisitos para Auto-Scaling Pleno}

Na implementação atual do projeto:
\begin{itemize}
    \item \textbf{Expansão de Servidores Backend:} Pode-se iniciar um Servidor 04 (porta 8004) e registrá-lo na inicialização do balanceador (\texttt{--servidores}). O balanceador redistribuirá instantaneamente a carga (passando de $33{,}3\%$ para $25\%$ por servidor), sem necessidade de reescrever ou recompilar o código;
    \item \textbf{Elasticidade de Consumidores:} Múltiplos processos consumidores (\texttt{consumidor.py}) podem se conectar concorrentemente à mesma fila do RabbitMQ (\textit{Competing Consumers Pattern}), e o broker dividirá a entrega das mensagens dinamicamente entre eles.
\end{itemize}

Para atingir a \textbf{escalabilidade automática plena}, o sistema exigiria três componentes de infraestrutura de nível de produção que ultrapassam o escopo pedagógico deste marco:
\begin{enumerate}
    \item \textbf{Sensor de Métricas e Gatilhos (\textit{Triggers}):} Um agente monitor que avaliasse a profundidade da fila do RabbitMQ e o tempo de resposta das requisições para disparar ordens de escala;
    \item \textbf{Orquestrador Dinâmico de Contêineres:} Uma ferramenta de orquestração (como Kubernetes HPA ou Docker Swarm) encarregada de clonar e iniciar novos nós automaticamente sob demanda;
    \item \textbf{Descoberta Dinâmica de Serviços (\textit{Service Discovery}):} Um registro centralizado em tempo real (como HashiCorp Consul ou Apache Zookeeper) que notificasse o balanceador da chegada do novo servidor, atualizando o algoritmo Round Robin sem intervenção humana ou reinicialização de processos.
\end{enumerate}

% =========================================================================
\section{Conclusão e Síntese Comparativa}
% =========================================================================

As modificações consolidadas elevaram o \textbf{Sistema Distribuído para Auxílio de Dados Hospitalares} de um modelo de protótipo síncrono inicial para uma arquitetura distribuída resiliente, desacoplada e aderente às melhores práticas de Engenharia de Sistemas Distribuídos. 

A Tabela~\ref{tab:comparativo_evolucao} sintetiza as transformações técnicas entre a versão original e o estado atual do projeto:

\begin{table}[H]
\centering
\small
\begin{tabularx}{\textwidth}{l X X}
\toprule
\textbf{Dimensão Técnica} & \textbf{Estado Inicial (Protótipo)} & \textbf{Estado Atual (Evolução com MOM e 6 VMs)} \\
\midrule
\textbf{Comunicação} & Síncrona via HTTP (Acoplada e bloqueante) & Assíncrona via protocolo AMQP 0-9-1 com desacoplamento temporal pleno \\
\addlinespace
\textbf{Bufferização de Tráfego} & Inexistente (picos de telemetria sobrecarregavam portas) & Fila durável RabbitMQ com persistência em disco (\texttt{delivery\_mode=2}) \\
\addlinespace
\textbf{Garantia de Entrega} & Frágil (quedas de servidor causavam perda de leituras) & Confirmação explícita (\texttt{basic\_ack}) e reenfileiramento (\texttt{basic\_nack}) \\
\addlinespace
\textbf{Controle de Fluxo} & Sem limite (risco de saturação de memória) & \texttt{prefetch\_count=1} garantindo entrega justa e prevenção de sobrecarga \\
\addlinespace
\textbf{Ambiente Virtual} & Servidores concorrendo na mesma VM & 6 Máquinas Virtuais totalmente isoladas com sistemas operacionais próprios \\
\addlinespace
\textbf{Segmentação de Rede} & Portas lógicas compartilhadas & Rede privada cabeada dedicada (\textit{Host-Only} \texttt{192.168.56.0/24}) \\
\addlinespace
\textbf{Observabilidade} & Visualização genérica de blocos & Painel SVG interativo com fronteiras reais de 6 VMs e conexão persistente SSE \\
\addlinespace
\textbf{Modelo de Escala} & Fixo e acoplado & Horizontal manual configurável com divisão proporcional de carga \\
\bottomrule
\end{tabularx}
\caption{Síntese comparativa da evolução arquitetural do sistema distribuído.}
\label{tab:comparativo_evolucao}
\end{table}

Em síntese, a integração do RabbitMQ como middleware de mensageria eliminou os riscos de perda silenciosa de sinais vitais críticos em momentos de sobrecarga. Concomitantemente, o desmembramento da infraestrutura em 6 Máquinas Virtuais independentes no \texttt{Vagrantfile} concedeu tolerância e isolamento real de falhas físicas e de sistema operacional, consolidando uma plataforma hospitalar robusta, auditável e didaticamente exemplar.

\vspace{0.4cm}
\begin{center}
    \rule{0.35\linewidth}{0.4pt}\\[4pt]
    {\footnotesize Departamento de Engenharia da Computação — UFMA • Outubro de 2026}
\end{center}

\end{document}
